\section{不可约表示与完全可约性}
\label{chap:schur-maschke}

一个给定的矩阵表示可能同时包含几种互不混合的对称类型。寻找不变子空间就是寻找这些块，但随意挑一个补空间未必仍对群作用不变。有限群的平均操作解决了这个困难：把任意内积在群上平均，便得到不变内积和正交不变补空间，这就是 Maschke 定理的构造内容。

分块以后，还要判断块之间允许哪些交织映射。Schur 引理说明不同不可约块之间没有非零交织，而一个复不可约块上与所有群元对易的算符只能是标量。群代数与正则表示把这两个结论放进同一个代数框架，也为下一章只用迹识别不可约重数做好准备。

\subsection{群代数与正则表示}
\label{sec:group-algebra-regular}


到目前为止，我们一直把表示理解成同态 $G\to \operatorname{GL}(V)$。另一种同样重要的观点是：把群元当作向量空间的一组基，从而把表示转化成一个代数模的理论。这样做的好处是，表示的直和、商、投影和中心元都会获得统一的代数解释。

\begin{definitionplain}[群代数]
设 $G$ 为有限群。其复群代数 $\mathbb C[G]$ 定义为所有形式和
\[
\sum_{g\in G} a_g\,g,
\qquad a_g\in\mathbb C,
\]
组成的复向量空间。以 $G$ 本身作为一组基，并把乘法按双线性延拓为
\[
\left(\sum_{g\in G}a_g g\right)
\left(\sum_{h\in G}b_h h\right)
=\sum_{g,h\in G}a_g b_h\,(gh).
\]
这样 $\mathbb C[G]$ 成为一个有限维结合代数，称为群 $G$ 的\textbf{群代数}。
\label{def:group-algebra}
\end{definitionplain}

群代数不是凭空出现的新对象。它只是把“对群元做线性组合”这一自然操作制度化。物理上常把这种线性组合看成某种“平均算子”或“投影算子”的来源；数学上，它使表示论变成模论，从而能够直接调用代数方法。

\begin{proposition}[表示与群代数模的对应]
设 $G$ 为有限群。给定一个表示 $\rho:G\to \operatorname{GL}(V)$，可在 $V$ 上定义 $\mathbb C[G]$ 的作用
\[
\left(\sum_{g\in G}a_g g\right)\cdot v
=\sum_{g\in G}a_g\,\rho(g)v.
\]
这使 $V$ 成为一个左 $\mathbb C[G]$-模。反过来，每个有限维左 $\mathbb C[G]$-模都唯一确定一个群表示。故有限维复表示与有限维左 $\mathbb C[G]$-模是一回事。
\label{prop:group-algebra-module}

\end{proposition}

\begin{proof}
先从表示到模。双线性由定义逐项继承。再检查乘法相容性：若
\[
\alpha=\sum_g a_g g,
\qquad
\beta=\sum_h b_h h,
\]
则
\begin{align*}
(\alpha\beta)\cdot v
&=\sum_{g,h}a_g b_h\,\rho(gh)v\\
&=\sum_{g,h}a_g b_h\,\rho(g)\rho(h)v\\
&=\alpha\cdot(\beta\cdot v).
\end{align*}
故确为左模。

反过来，若 $V$ 是左 $\mathbb C[G]$-模，对每个 $g\in G$ 定义 $\rho(g)v=g\cdot v$。由于群元在群代数中可乘，便有
\[
\rho(gh)v=(gh)\cdot v=g\cdot(h\cdot v)=\rho(g)\rho(h)v,
\]
因此 $\rho$ 是表示。两边构造互逆，命题得证。
\end{proof}

在所有表示中，最能反映群自身结构的是正则表示。它直接让群通过左乘或右乘作用在"以群元为基的空间"上，因此与\chapref{chap:group-basics}的重排定理密切相关。

\begin{definitionplain}[左正则表示]
设 $V_{\mathrm{reg}}$ 为以 $\{e_g\mid g\in G\}$ 为基的 $|G|$ 维复向量空间。定义
\[
L(h)e_g=e_{hg},
\qquad h,g\in G.
\]
其线性延拓给出表示
\[
L:G\to \operatorname{GL}(V_{\mathrm{reg}}),
\]
称为 $G$ 的\textbf{左正则表示}。
\label{def:left-regular-representation}
\end{definitionplain}

这个定义看似只是把群元“重命名”为基向量，再让群用左乘去重新排列它们。但它的意义非常大：由于每个基向量都对应一个群元，正则表示把整个群的信息都装进了一个表示空间里。很多一般定理的常数因子，最终都能在这个表示上通过简单计数求出来。

\begin{proposition}[左正则表示确为表示]
映射 $L:G\to\operatorname{GL}(V_{\mathrm{reg}})$ 满足
\[
L(h_1h_2)=L(h_1)L(h_2),
\qquad
L(e)=I.
\]
因此它是 $G$ 的表示。并且对每个固定的 $h\in G$，矩阵 $L(h)$ 在基 $\{e_g\}$ 下是一个置换矩阵。
\label{prop:left-regular-is-representation}

\end{proposition}

\begin{proof}
对任意基向量 $e_g$，有
\[
L(h_1h_2)e_g=e_{h_1h_2g}=L(h_1)e_{h_2g}=L(h_1)L(h_2)e_g,
\]
故 $L(h_1h_2)=L(h_1)L(h_2)$。又有 $L(e)e_g=e_{eg}=e_g$，所以 $L(e)=I$。当 $h$ 固定时，映射 $g\mapsto hg$ 是集合 $G$ 到自身的双射，因此 $L(h)$ 只是把基向量做一个重排，故其矩阵为置换矩阵。
\end{proof}

左正则表示在基 $\{e_g\}$ 上的重排结构如图 \figref{fig:left-regular-action} 所示。
\begin{figure}[htbp]
\centering
\begin{tikzfig}
  \node[lecturebox] (basis) at (0,0) {$V_{\mathrm{reg}}$ 的基\\$\{e_g\}_{g\in G}$};
  \node[lecturebox] (left) at (7.2,0) {固定 $h\in G$\\$L(h)e_g=e_{hg}$};
  \draw[lecturearrow] (basis) -- node[above,lecturelabel] {左乘重排} (left);
  \node[lecturelabel,align=left] at (3.6,-1.6) {由于 $g\mapsto hg$ 是双射，\\$L(h)$ 在这组基下总是置换矩阵。};
\end{tikzfig}
\caption{左正则表示的最核心结构：固定一个 $h$，它只是在全部群元标签之间做左乘重排。}
\label{fig:left-regular-action}
\end{figure}

\medskip
\noindent\emph{为什么正则表示包含整个群的信息。}
若 $G$ 自己通过左乘作用在集合 $G$ 上，那么由 \Cref{ex:permutation-representation}，这会诱导出一个置换表示。此时作用集合本身就是群，而诱导出的置换表示恰好就是左正则表示。因此正则表示并不是凭空新增的例子，而是“群作用 $\Rightarrow$ 置换表示”这一一般机制在最天然作用上的具体实现。
\par\medskip

正则表示可以在 $C_3$ 上看得很具体。取基 $\{e_e,e_a,e_{a^2}\}$，左乘 $a$ 的作用是
\[
e_e\mapsto e_a,\qquad e_a\mapsto e_{a^2},\qquad e_{a^2}\mapsto e_e,
\]
所以
\[
L(a)=
\begin{pmatrix}
0&0&1\\
1&0&0\\
0&1&0
\end{pmatrix}.
\]
这个三维置换矩阵的本征值是 $1,\omega,\omega^2$，其中 $\omega=\ee^{2\pi\ii/3}$。对应三个一维本征子空间，正好就是 $C_3$ 的三个不可约复表示。因此“正则表示包含所有不可约表示”在这个例子里就是普通矩阵对角化。

\subsection{有限群表示理论的基本结构}
\label{sec:finite-group-structure}


上一节已经看到：若表示是酉的，那么不变子空间的正交补也不变，因此表示可以分裂成直接和。结合 \Cref{thm:unitarization} 立刻得到有限群表示论的第一主定理，即 Maschke 定理。

\begin{theorem}[Maschke 定理]
设 $G$ 为有限群，$V$ 为其有限维复表示。若 $W\subseteq V$ 是不变子空间，则存在另一个不变子空间 $W'$，使得
\[
V=W\oplus W'.
\]
因此任意有限维复表示都可分解为有限个不可约表示的直接和。
\label{thm:maschke}

\end{theorem}

\begin{proof}
由 \Cref{thm:unitarization}，$V$ 上可选取某个厄米内积使表示酉化。于是根据 \Cref{prop:orthogonal-complement-invariant}，$W^{\perp}$ 也是不变子空间，并且由正交分解定理
\[
V=W\oplus W^{\perp}.
\]
这就给出了所需的不变互补。

要说明“任意表示分解为不可约表示直接和”，只需对维数作有限下降。若 $V$ 不可约则已完成；若可约，则存在非平凡不变子空间 $W$，由前一步得到
\[
V=W\oplus W'.
\]
若 $W$ 或 $W'$ 仍可约，就继续分裂。由于每一步都把表示分裂成维数更小的子表示，而总维数有限，故该过程必在有限步后终止，终点就是若干不可约表示的直接和。
\end{proof}

上一证明利用“有限群表示可以酉化”，推导很短。另一个更值得掌握的证明直接展示了\emph{群平均怎样把一个任意投影修正成群不变投影}。先只用普通线性代数选一个补空间 $V=W\oplus W_0$，并令 $P_0:V\to W$ 是沿 $W_0$ 投到 $W$ 的投影。这个 $P_0$ 一般不与群作用对易，所以 $W_0$ 一般也不是不变子空间。对它做群平均：
\begin{equation}
P=\frac1{|G|}\sum_{g\in G}\rho(g)P_0\rho(g)^{-1}.
\label{eq:maschke-averaged-projector}
\end{equation}
先看 $P$ 是否仍然投到 $W$。因为 $W$ 本身对 $G$ 不变，任取 $v\in V$，$P_0\rho(g)^{-1}v\in W$，再作用 $\rho(g)$ 仍在 $W$，所以每一项都落在 $W$，从而 $Pv\in W$。

更关键的是，对 $w\in W$，有 $\rho(g)^{-1}w\in W$，而 $P_0$ 在 $W$ 上就是恒等算符，于是每个求和项都满足
\[
\rho(g)P_0\rho(g)^{-1}w
=\rho(g)\rho(g)^{-1}w=w.
\]
因此 $Pw=w$。也就是说 $P$ 的像正是 $W$，并且 $P$ 在自己的像上等于恒等算符，所以
\[
P^2=P.
\]

最后验证这个平均后的投影真的尊重群作用。对任意 $h\in G$，
\begin{align*}
\rho(h)P\rho(h)^{-1}
&=\frac1{|G|}\sum_{g\in G}
\rho(hg)P_0\rho(hg)^{-1}\\
&=\frac1{|G|}\sum_{k\in G}
\rho(k)P_0\rho(k)^{-1}=P,
\end{align*}
第二行只是把求和变量重命名为 $k=hg$；左乘 $h$ 会把有限群的全部元素重新排列一次，不会改变求和。因此 $P\rho(h)=\rho(h)P$。

现在取
\[
W'=\operatorname{Ker} P.
\]
若 $v\in W'$，则
\[
P\rho(h)v=\rho(h)Pv=0,
\]
所以 $\rho(h)v\in W'$，即 $W'$ 也是不变子空间。又因为 $P$ 是投影到 $W$，线性代数给出
\[
V=\operatorname{Im} P\oplus\operatorname{Ker} P=W\oplus W'.
\]
这就再次证明了 Maschke 定理。这个版本以后非常有用，因为“先随便构造一个对象，再对整个有限群平均，使它自动具有对称性”会在内积、投影算符和选择定则中反复出现。

这里“完全可约”比“可约”强得多。一个表示可约，只说明存在某个不变子空间 $W$；但要写成直和 $V=W\oplus W'$，还必须找到一个同样不变的补空间 $W'$。一般线性代数当然总能找到某个补空间，但它未必被群作用保持。

有限群的关键优势是可以先酉化表示。酉表示中 $W^\perp$ 自动不变，于是可以取 $W'=W^\perp$。如果 $W$ 或 $W'$ 还可约，就继续分。维数每次严格下降，所以有限维情况下不可能无限进行。最终得到
\[
V\simeq\bigoplus_\alpha m_\alpha V_\alpha.
\]
后面所有“用特征标直接求重数”的公式都默认了这一步已经成立。


Maschke 定理是有限群表示论的分水岭。它说明在复数域上，有限群表示不存在“粘连不散”的病态块，所有表示都可完全拆开。后面的特征标理论之所以有效，根本原因就在于这种完全可约性。

物理上，Maschke 定理保证了具有有限群对称性的哈密顿量 $\hat H$ 的能级简并可以被群论完全分类。设 $\hat H$ 与群 $G$ 的所有表示算符 $U(g)$ 交换，则每个能级 $E$ 的本征空间 $\mathcal E_E$ 是 $G$ 的表示空间。由 Maschke 定理，$\mathcal E_E$ 可以完全分解为不可约表示的直接和：$\mathcal E_E\simeq\bigoplus_\alpha m_\alpha V_\alpha$。每个不可约表示 $V_\alpha$ 贡献 $d_\alpha$ 维简并，因此能级 $E$ 的简并度
\[
\dim\mathcal E_E=\sum_\alpha m_\alpha d_\alpha
\]
完全由群论决定。如果没有 Maschke 定理——例如对某些无限群或实数域上的表示——可能存在"粘连不散"的简并块，群论无法将其拆开，简并度就不能仅由对称性预测。$\mathrm{NH_3}$ 的振动模式 $2A_1\oplus 2E$（\eqnrefs{eq:nh3-vib-decomposition}）之所以能从特征标表直接读出简并结构，正是 Maschke 定理在背后保证了振动表示可以完全分解。

\begin{definitionplain}[重数与同类分量]
设表示 $V$ 分解为不可约表示的直接和
\[
V\simeq \bigoplus_{\alpha} m_\alpha V_\alpha,
\]
其中 $V_\alpha$ 取遍一组两两不等价的不可约表示，$m_\alpha\in\mathbb N\cup\{0\}$。则 $m_\alpha$ 称为不可约表示 $V_\alpha$ 在 $V$ 中出现的\textbf{重数}。固定某个 $\alpha$ 时，
\[
V[\alpha]\simeq m_\alpha V_\alpha
\]
称为与 $V_\alpha$ 对应的\textbf{同类分量}（等型分量）。
\label{def:multiplicity-isotypic}
\end{definitionplain}

完全可约性保证了重数的良定义性：虽然具体直和分解可能因基底与互补的选取而不同，但每个不可约等价类出现多少次是唯一不变的。这一点在特征标部分会由正交关系给出一个显式公式。

\begin{theorem}[表示分解的唯一性（Krull--Schmidt 型）]
\label{thm:decomposition-uniqueness}
设有限群 $G$ 的表示 $V$ 有两种不可约分解
\[
V\cong\bigoplus_\alpha m_\alpha V_\alpha
\cong\bigoplus_\alpha m'_\alpha V_\alpha,
\]
其中 $V_\alpha$ 取遍不等价不可约表示。则 $m_\alpha=m'_\alpha$ 对所有 $\alpha$ 成立。

\end{theorem}

\begin{proof}
对每个 $\beta$，计算 $\operatorname{Hom}_G(V_\beta,V)$（$G$-等变线性映射的空间）。由 Schur 引理
\[
\operatorname{Hom}_G(V_\beta,V_\alpha)=\begin{cases}\CC,&\alpha=\beta,\\0,&\alpha\neq\beta.\end{cases}
\]
因此
\[
\operatorname{Hom}_G(V_\beta,V)\cong\bigoplus_\alpha m_\alpha\operatorname{Hom}_G(V_\beta,V_\alpha)=m_\alpha\CC,
\]
故 $\dim\operatorname{Hom}_G(V_\beta,V)=m_\beta$。同理 $\dim\operatorname{Hom}_G(V_\beta,V)=m'_\beta$，故 $m_\beta=m'_\beta$。
\end{proof}

\begin{proposition}[投影算符的显式公式]
\label{prop:projection-operator}
设 $V=\bigoplus_\alpha m_\alpha V_\alpha$ 是 $G$ 的表示 $V$ 的不可约分解。对每个不可约表示 $\alpha$，投影到 $\alpha$-同类分量 $V[\alpha]\cong m_\alpha V_\alpha$ 的 $G$-等变投影算符为
\begin{equation}
P_\alpha=\frac{d_\alpha}{|G|}\sum_{g\in G}\overline{\chi_\alpha(g)}\,\rho(g).
\label{eq:projection-operator}
\end{equation}
它满足
\[
P_\alpha^2=P_\alpha,\qquad P_\alpha P_\beta=0\ (\alpha\neq\beta),\qquad\sum_\alpha P_\alpha=I_V.
\]

\end{proposition}

\begin{proof}
对任意 $h\in G$，共轭只是在求和中重排群元，而特征标在共轭类上常值，因此
\[
\rho(h)P_\alpha\rho(h)^{-1}=\frac{d_\alpha}{|G|}\sum_g\overline{\chi_\alpha(g)}\rho(hgh^{-1})
=\frac{d_\alpha}{|G|}\sum_{g'}\overline{\chi_\alpha(g')}\rho(g')
=P_\alpha,
\]
故 $P_\alpha$ 与表示的作用交换。

将它限制到不可约分量 $V_\beta$ 上。Schur 引理给出
$P_\alpha|_{V_\beta}=\lambda_\beta I_{V_\beta}$，而取迹可得
\[
\Tr_{V_\beta}P_\alpha=\frac{d_\alpha}{|G|}\sum_g\overline{\chi_\alpha(g)}\chi_\beta(g)=d_\alpha\langle\chi_\alpha,\chi_\beta\rangle=d_\alpha\delta_{\alpha\beta}.
\]
于是 $\lambda_\beta d_\beta=d_\alpha\delta_{\alpha\beta}$，即
$\lambda_\beta=\delta_{\alpha\beta}$。所以 $P_\alpha$ 在 $V[\alpha]$ 上为恒等、在其余同类分量上为零；幂等性、相互正交性与 $\sum_\alpha P_\alpha=I$ 随即成立。
\end{proof}

\begin{exampleplain}[投影算符在对称性适配线性组合构造中的应用]
\label{ex:projection-salc}
在分子物理中，投影算符 \eqnrefs{eq:projection-operator} 直接用于从原子轨道构造对称适配线性组合（对称性适配线性组合）。以 $\mathrm{NH_3}$（$C_{3v}\cong D_3$）为例，三个 H 原子的 $1s$ 轨道 $\{\phi_1,\phi_2,\phi_3\}$ 张成三维空间，承载 $C_{3v}$ 的表示 $A_1\oplus E$。

投影到 $A_1$（平凡表示，$d=1$，$\chi_{A_1}=(1,1,1)$）：
\[
P_{A_1}=\frac{1}{6}\bigl(E+C_3+C_3^2+\sigma_v+\sigma_v'+\sigma_v''\bigr).
\]
作用在 $\phi_1$ 上：
\[
P_{A_1}\phi_1=\frac{1}{6}(\phi_1+\phi_2+\phi_3+\phi_1+\phi_3+\phi_2)=\frac{1}{3}(\phi_1+\phi_2+\phi_3),
\]
归一化后给出 $A_1$ 对称的对称性适配线性组合 $\psi_{A_1}=\frac{1}{\sqrt3}(\phi_1+\phi_2+\phi_3)$。

投影到 $E$（二维表示，$d=2$，$\chi_E=(2,0,-1)$）：
\[
P_E=\frac{d_E}{|G|}\sum_{g\in C_{3v}}\chi_E(g)^*\,g
  =\frac{1}{3}\bigl(2E-C_3-C_3^2\bigr).
\]
作用在 $\phi_1$ 上给出 $E$ 表示的一个基矢。这就是\chapref{chap:group-theory-quantum} 中对称性适配线性组合构造的群论基础——投影算符公式把"从对称性提取适配基"变成了一个机械计算。
\end{exampleplain}

前面的定理仍然是在研究"表示空间里的向量"。现在把视角换一下：对固定指标 $i,j$，矩阵元
\[
g\longmapsto D^{(\alpha)}_{ij}(g)
\]
本身是定义在有限集合 $G$ 上的复值函数。所有这样的函数构成一个 $|G|$ 维复向量空间，我们可以像对普通向量一样定义内积
\[
\langle f,h\rangle_G=\frac1{|G|}\sum_{g\in G}\overline{f(g)}h(g).
\]
所谓“大正交关系”，就是说不同不可约表示的矩阵元在这个函数空间里恰好像互相垂直的向量。这样一来，表示论的分解问题就能被转化成我们熟悉的“正交展开”。

有限群表示论中最有力量的计算工具之一是矩阵元的正交性。它的推导恰好展示了 Schur 引理为何如此关键：先构造一个平均化的交织算子，再由不可约性迫使该算子只能是标量乘恒等算子。

\begin{theorem}[大正交关系]
设 $G$ 为有限群，$D^{(\alpha)}$ 与 $D^{(\beta)}$ 是两个不可约酉表示，其维数分别为 $d_\alpha,d_\beta$。则对任意矩阵元指标 $i,j,k,\ell$，有
\[
\sum_{g\in G} \overline{D^{(\alpha)}_{ij}(g)}\,D^{(\beta)}_{k\ell}(g)
=\frac{|G|}{d_\alpha}\,\delta_{\alpha\beta}\delta_{ik}\delta_{j\ell}.
\]
\label{thm:great-orthogonality}

\end{theorem}

\begin{proof}
先固定两个不可约酉表示 $D^{(\alpha)}$ 与 $D^{(\beta)}$。任取线性映射 $A:V_\beta\to V_\alpha$，构造群平均
\[
M_A=\sum_{g\in G}D^{(\alpha)}(g)A D^{(\beta)}(g)^{-1}.
\]
这个平均的目的，是把任意线性映射 $A$ 投影到与群作用相容的交织算子空间。对任意 $h\in G$，直接同时在左右作用群元：
\begin{align*}
D^{(\alpha)}(h)M_A D^{(\beta)}(h)^{-1}
&=\sum_{g\in G}
D^{(\alpha)}(hg)A
D^{(\beta)}(g)^{-1}D^{(\beta)}(h)^{-1}\\
&=\sum_{g\in G}
D^{(\alpha)}(hg)A D^{(\beta)}(hg)^{-1}.
\end{align*}
令 $k=hg$。左乘 $h$ 只是把有限集合 $G$ 重排一次，所以 $k$ 仍遍历 $G$，从而
\[
D^{(\alpha)}(h)M_A D^{(\beta)}(h)^{-1}=M_A.
\]
等价地，
\[
D^{(\alpha)}(h)M_A=M_A D^{(\beta)}(h),
\]
所以 $M_A$ 确实是从 $V_\beta$ 到 $V_\alpha$ 的交织算子。

若 $\alpha\neq\beta$，Schur 引理给出 $M_A=0$。若 $\alpha=\beta$，则 $M_A=c(A)I$。此时取迹，利用迹的循环性，
\begin{align*}
\Tr M_A
&=\sum_{g\in G}\Tr\bigl(D^{(\alpha)}(g)A D^{(\alpha)}(g)^{-1}\bigr)\\
&=|G|\Tr A.
\end{align*}
另一方面 $\Tr M_A=d_\alpha c(A)$，故
\[
c(A)=\frac{|G|}{d_\alpha}\Tr A.
\]
因此对任意 $A$，
\begin{equation}
M_A=
\begin{cases}
0,&\alpha\neq\beta,\\[1mm]
\dfrac{|G|}{d_\alpha}\Tr(A)I,&\alpha=\beta.
\end{cases}
\label{eq:averaged-intertwiner}
\end{equation}
现在选取正交归一基，并令 $A=E_{j\ell}$ 为矩阵单位：它把 $V_\beta$ 的第 $\ell$ 个基矢送到 $V_\alpha$ 的第 $j$ 个基矢。于是 $M_A$ 的 $(i,k)$ 矩阵元为
\begin{align*}
(M_A)_{ik}
&=\sum_{g\in G}
D^{(\alpha)}_{ij}(g)
\bigl[D^{(\beta)}(g)^{-1}\bigr]_{\ell k}\\
&=\sum_{g\in G}
D^{(\alpha)}_{ij}(g)
\overline{D^{(\beta)}_{k\ell}(g)},
\end{align*}
其中第二行使用了酉性 $D(g)^{-1}=D(g)^\dagger$。另一方面，由 \eqnrefs{eq:averaged-intertwiner}，
\[
(M_A)_{ik}
=\frac{|G|}{d_\alpha}
\delta_{\alpha\beta}\delta_{j\ell}\delta_{ik}.
\]
两式比较后取复共轭，就得到定理所述
\[
\sum_{g\in G}
\overline{D^{(\alpha)}_{ij}(g)}D^{(\beta)}_{k\ell}(g)
=\frac{|G|}{d_\alpha}
\delta_{\alpha\beta}\delta_{ik}\delta_{j\ell}.
\]
\end{proof}
大正交关系里的四个指标很容易把初学者淹没，可以先只看它在说什么。对每个不可约表示 $\alpha$，每个矩阵元 $D^{(\alpha)}_{ij}$ 都是一条“定义在群上的函数”。定理说这些函数在离散内积下彼此正交，归一化长度为 $|G|/d_\alpha$：
\[
\left\langle D^{(\alpha)}_{ij},D^{(\beta)}_{k\ell}\right\rangle_G
=\frac1{|G|}\sum_g \overline{D^{(\alpha)}_{ij}(g)}D^{(\beta)}_{k\ell}(g)
=\frac1{d_\alpha}\delta_{\alpha\beta}\delta_{ik}\delta_{j\ell}.
\]
三个 Kronecker \(\delta\) 分别表达三层正交：不同不可约表示正交；同一个不可约表示中不同行指标正交；不同列指标也正交。
证明为什么会得到这个形式？群平均把任意矩阵 $A$ 变成一个交织算子；Schur 引理再把交织算子压缩成 $cI$。因此“大正交关系 = 群平均 + Schur 引理 + 取矩阵单位”。把这条因果链记住，比死记四指标公式更重要。
这一公式的本质应当理解为：不同不可约表示的矩阵元在群上的离散求和内积下彼此正交，而同一不可约表示内部的不同矩阵元也两两正交。这种“比特征标更强”的正交性在研究投影算符与选择定则时非常有用；特征标正交关系只是它在对角元求和后的简化版本。
